\documentclass[11pt,a4paper]{article}
\usepackage[margin=27mm]{geometry}
\usepackage{amsmath,amssymb,amsthm,hyperref}
\newtheorem{lemma}{Lemma}\newtheorem{theorem}[lemma]{Theorem}
\newcommand{\Sz}{\operatorname{Sz}}\newcommand{\F}{\mathbb F}
\title{Extensions of the split affine residues}\author{Human Compression project}\date{}
\begin{document}\maketitle
\begin{abstract}
We determine every rank-four residually weakly primitive extension of
the specified residues R7 and R9 in the fixed higher-rank Suzuki paper.
The proof uses the actual prescribed intersections and does not assume
that paper's classification of all rank-three residues. The main step
is a subfield containment lemma for an affine subgroup whose scalar
complement generates its coefficient field.
\end{abstract}

\section{Conventions and structural premises}
Write $S_f=\Sz(2^f)$ for the standard subfield group in $S_N$, where
$f\mid N$ and both exponents are odd. All simple subfield groups below
have $f\ge3$. Write $H_m$ for the subgroup of order $m\mid2^N-1$ in
the standard split torus. Its Weyl involution $w$ inverts it. Write
$Z_{2^f}$ for the involution centre of the standard Sylow two-subgroup
$Q_{2^f}$. In particular, $Z_{2^f}$ has the additive coordinates of
$\F_{2^f}$. Put $D_{2m}=H_m\langle w\rangle$.

We use the classical Suzuki matrix model, its subgroup classification,
and conjugacy of subfield groups and split tori, as in the accepted
preliminary papers. Specifically, a proper nonsoluble subgroup is a
simple subfield group. The other maximal subgroup types are a point
stabilizer, a split torus normalizer, and one of the two nonsplit torus
normalizers. A nonidentity odd-order element lies in a unique torus;
$N_{S_N}(H_m)$ is the full split normalizer for $m>1$. Simple subfield
groups are self-normalizing. A point stabilizer has shape
$Q_q:H_{q-1}$, where $|Z_q|=|Q_q/Z_q|=q$ and $[Q_q,Q_q]=Z_q$.
The torus acts on both additive quotients by scalar powers coprime to
$q-1$. A nonidentity involution fixes a unique ovoid point. These are
explicit structural premises, not consequences of the classification
being corrected.

The source's residual weak primitivity condition says that, for every
proper subset $J$ of the type set, some remaining intersection
$G_{J\cup\{i\}}$ is maximal proper in $G_J$. It implies that every
nested intersection is strict. Indeed, if $G_J\le G_i$ for $i\notin J$,
then $G_{I\setminus\{i\}}\le G_J\le G_i$, contradicting maximal
proper inclusion at $I\setminus\{i\}$. We use this consequence
explicitly; no geometric flag transitivity is assumed.

\begin{lemma}\label{simple}
A fixed embedded simple subfield group $S_e$ is contained in at most
one subfield group of any prescribed degree $f$.
\end{lemma}
\begin{proof}
Suppose a conjugate of $S_f$ contains $S_e$. Transporting to $S_f$ and
using conjugacy of its subfield groups, adjust the transporter so that
it normalizes $S_e$. Its ambient normalizer is $S_e$, which is already
contained in the standard $S_f$. Thus the conjugate is $S_f$ itself.
\end{proof}

\section{A field-generating affine complement}
\begin{lemma}\label{affine}
Let $s=2^f\ge8$, let $f\mid N$, and let $m\mid s-1$ such that $H_m$
generates $\F_s$ as a field over $\F_2$. Any simple subfield group
containing either $D_{2m}$ or $Z_s:H_m$ contains the standard $S_f$.
\end{lemma}
\begin{proof}
Let $T$ have degree $v$. First $H_m$ is split in $T$. Indeed, the order
of $2$ modulo $m$ is $f$, an odd number. If $H_m$ were in a nonsplit
torus, then $m\mid 2^{2v}+1$, so this order would divide $4v$ but
not $2v$, which is impossible for an odd order. Consequently
$m\mid2^v-1$, and field generation gives $f\mid v$.

All subgroups of order $m$ in split tori of $S_v$ are conjugate.
If $T=gS_vg^{-1}$ contains $H_m$, adjust $g$ by an element of $S_v$
so that it normalizes $H_m$. The ambient normalizer is the split
dihedral group. Its reflection normalizes $S_v$, so we may take
$g=h\in H_{q-1}$, where $q=2^N$.

If $w\in T$, then $h^{-1}wh=h^{-2}w\in S_v$, whence $h^2\in H_{2^v-1}$.
All torus orders are odd, so $h\in H_{2^v-1}$ and $T=S_v$.
This treats $D_{2m}$.

For the affine case, $T\cap Z_q=hZ_{2^v}h^{-1}$ has additive coordinate
set $\lambda\F_{2^v}$, where $\lambda$ is the scalar by which $h$ acts
on $Z_q$. Containment of $Z_s$ implies
$1\in\F_s\subseteq\lambda\F_{2^v}$, hence $\lambda\in\F_{2^v}$.
The central scalar exponent is coprime to $q-1$. Taking its inverse
power shows that $h\in H_{2^v-1}$. Again $T=S_v$, which contains $S_f$.
\end{proof}

We also use the following acquired arithmetic fact. If $t=2^e\ge8$,
$a$ is an odd prime, and $n\mid(t^a-1)/(t-1)$ is prime, then
$H_{(t-1)n}$ generates $\F_{t^a}$. If $n\nmid t-1$, the order of
$t$ modulo $n$ is $a$. If $n\mid t-1$, then $n=a$ and the valuation
of $n$ in $t^a-1$ is exactly one larger than in $t-1$. In either case
no proper intermediate field contains this cyclic subgroup. The full
valuation proof, including overlap with $t-1$, is in the accepted
paper \emph{Scalar fields and split torus families}.

\section{All extensions of the prescribed residue}
\begin{theorem}\label{extension}
Let a rank-four residually weakly primitive coset pregeometry in
$G=\Sz(q)$ have $G_0=\Sz(s)$ maximal in $G$. Suppose its prescribed
$G_0$-residue is R7 or R9. In explicit coordinates put
\[
 t=2^e\ge8,\quad s=t^a,\quad b=t-1,\quad m_0=bn,
 \qquad n\mid(s-1)/b\text{ prime},
\]
where $a$ is prime. The prescribed lattice is
\[
 G_{01}=\Sz(t),\quad G_{02}=D_{2m_0},\quad G_{03}=Z_s:H_{m_0},
\]
\[
 G_{012}=D_{2b},\quad G_{013}=Z_t:H_b,\quad
 G_{023}=H_{m_0},\quad G_{0123}=H_b.
\]
After one simultaneous conjugation, every extension has
\[
 q=t^{ac},\quad G_1=\Sz(v),\quad v=t^c,
 \qquad a\ne c\text{ primes},
\]
and, for a prime $n'\ne n$ dividing $(v-1)/b$,
\[
 G_2=D_{2bnn'},\qquad G_3=Z_q:H_{bnn'}.
\]
All remaining intersections are
\[
 G_{12}=D_{2bn'},\quad G_{13}=Z_v:H_{bn'},\quad
 G_{23}=H_{bnn'},\quad G_{123}=H_{bn'}.
\]
\end{theorem}
\begin{proof}
We first justify the coordinates in the statement. Conjugate $G_0$
and its simple subfield $G_{01}$ to the standard copies. The full
split normalizer $G_{012}$ in $G_{01}$ may then be made
$H_b\langle w\rangle$. The odd rotation subgroup of $G_{02}$ lies
in the unique ambient split torus containing $H_b$, so
$G_{02}=H_{m_0}\langle w\rangle$. The group $G_{03}$ is an affine
subgroup of $G_0$ with elementary abelian radical of order $s$.
Its radical is the centre of a Sylow two-subgroup of $G_0$.
Since $G_{023}=H_{m_0}$, its fixed point is one of the two torus axis
points. Interchanging them by $w$ if necessary gives
$G_{03}=Z_s:H_{m_0}$. This also gives the displayed coordinates
for $G_{013}$ and all deeper prescribed intersections.

Since $G_0$ is maximal, $q=s^c$ for a prime $c$. The proper group
$G_1$ contains the simple $\Sz(t)$, hence is itself a simple subfield
group. Lemma~\ref{simple} makes it standard. Its degree is a proper
divisor of $eac$, contains $e$, and has gcd $e$ with $ea$ because
$G_0\cap G_1=\Sz(t)$. Therefore its degree is $ec$, with $a\ne c$.
Thus $G_1=\Sz(v)$, $v=t^c$.

The scalar group $H_{m_0}$ generates $\F_s$. Lemma~\ref{affine}
and maximality of $S_{ea}$ show that no proper subfield overgroup
of $G_{02}$ or $G_{03}$ can lie outside $G_0$. Strictness excludes
containment of $G_2$ or $G_3$ in $G_0$. A point stabilizer cannot
contain $D_{2m_0}$: the odd torus fixes its two axis points and the
reflection interchanges them. Nonsplit torus normalizers are also
excluded by their odd orders. Hence
\[
 G_2=D_{2M},\quad M\mid q-1,\qquad
 m_0=\gcd(M,s-1),\quad m_1:=\gcd(M,v-1)>b.
\]
Here $G_{12}=D_{2m_1}$ and $\gcd(m_0,m_1)=b$.
The elementary abelian subgroup $Z_s$ has order at least eight,
excluding both sorts of torus normalizer as overgroups of $G_3$.
Its fixed point is unique. Thus $G_3$ lies in the standard point
stabilizer $B_q=Q_q:H_{q-1}$.

Write $G_{23}=H_h$. This is cyclic, since the reflection coset of
$G_2$ is outside $B_q$. Strictness and the prescribed intersection
with $G_0$ give $h>m_0$ and $\gcd(h,s-1)=m_0$. The field generated
by $H_h$ contains $\F_s$ and is not $\F_s$; since $[\F_q:\F_s]=c$
is prime, it is $\F_q$.

Put $N=G_3\cap Q_q$. Its intersection with $Z_q$ contains $Z_s$.
The scalar action of $H_h$ is irreducible on $Z_q$, so
$N\cap Z_q=Z_q$. Its image on $Q_q/Z_q$ is also invariant under
this irreducible scalar action. Thus $N=Z_q$ or $N=Q_q$.
The second case would make $G_0\cap G_3$ contain $Q_s$, contradicting
the prescribed central affine group. Therefore $N=Z_q$.

The quotient $G_3/Z_q$ maps injectively into the cyclic complement
of $B_q/Z_q$, and is therefore abelian. It contains the nontrivial
scalar subgroup $H_h$. In the affine quotient $Q_q/Z_q:H_{q-1}$,
a translation commuting with a nonidentity scalar must be zero.
Every element of $G_3/Z_q$ therefore has zero translation coordinate.
Consequently
\[
 G_3=Z_q:H_d,\qquad h=\gcd(M,d),\qquad m_0\mid h\mid d.
\]
Put $d_1=\gcd(d,v-1)$. Both $m_0<d$ and $d_1<d$:
the former follows from $h>m_0$, and the latter from field generation
by $H_{m_0}$, since $\F_s\nsubseteq\F_v$.
The intersections $G_{03}$ and $G_{13}$ have proper intermediate
overgroups $Z_q:H_{m_0}$ and $Z_q:H_{d_1}$ in $G_3$.
If $h<d$, then $G_{23}=H_h<H_d<G_3$ as well. This would violate
residual weak primitivity at $G_3$. Hence $h=d$, so $d\mid M$.

We claim $d=M$. Suppose $d<M$. Then $H_d$ is not maximal in
$D_{2M}$, because $H_d<H_M<D_{2M}$. Residual weak primitivity
at $G_2$ makes either $D_{2m_0}$ or $D_{2m_1}$ maximal.
For odd dihedral groups this means $M/m_0$ or $M/m_1$ is prime:
all overgroups of a fixed reflection are obtained by enlarging its
cyclic rotation subgroup.
If $M/m_0$ is prime, the strict divisor chain $m_0<d\mid M$ gives
$d=M$, a contradiction. Otherwise $M/m_1$ is prime. Since
$m_0/\gcd(m_0,m_1)=n$ divides it, we have $M=nm_1$.
As $m_0\mid d<M=\operatorname{lcm}(m_0,m_1)$, one has
$\gcd(d,m_1)<m_1$. In the residue at $G_{12}=D_{2m_1}$, the
intersection with $G_3$ is then a proper subgroup of $H_{m_1}$
and cannot be maximal. The intersection with $G_0$, namely $D_{2b}$,
must be maximal. Thus $m_1/b$ is prime. But
$M/m_0=m_1/b$, again forcing $d=M$ from $m_0<d\mid M$.
This contradiction proves the claim.

Now $G_{123}=H_{m_1}$. The final residual condition makes
$H_b$ maximal in it, so $m_1=bn'$ for a prime $n'$.
The equality $\gcd(m_0,m_1)=b$ gives $n'\ne n$.
At $G_{23}=H_M$, at least one of $H_{m_0}$ and $H_{m_1}$ is
maximal. The prime index must respectively be $n'$ or $n$, because
the other subgroup is contained in $H_M$. Therefore $M=bnn'$.
The coordinate intersections now give every asserted formula.
\end{proof}

\section{Exact source scope and attribution}
The source is D.~Leemans, \emph{The residually weakly primitive
pre-geometries of the Suzuki simple groups}, Note di Matematica
20(1) (2000/2001), 1--20, fixed twenty-page edition.
R7 is discussed on p.~13 and again in Section~3.3.2; R9 is on p.~14.
In R7, $e$ is prime; R9 allows composite $e$. The proof above works
in both cases. It supplies the conditional necessity missing from a
proof of the canonical Constructions 5 and 7 alone. It proves neither
the source's false assertion that these and the other listed residues
exhaust all cases, nor a classification of unrestricted relative
positions outside the prescribed residue.

There are also literal errors in the R7 paragraph. Its final affine
group is called $G_4$ although the types are $0,1,2,3$; it is $G_3$.
Its printed condition $n'\mid(q-1)/(s'-1)$ does not express the
required additional scalar factor. The correct condition is
$n'\mid(s'-1)/(t-1)$, with $s'=v$, as printed in the R9 paragraph.
The theorem gives the corrected full lattice without using either
of these inconsistent expressions.

The classical ambient structure and conjugacy facts are borrowed as
stated. The scalar valuation argument is reused from the accepted
split-family proof. The coordinate containment lemma and the necessity
argument above are the present contribution. They repair a gap in our
first whole-source claim ledger, which confused a refuted global
classification with these separate conditional assertions. No new
finite search, Lean declaration, principal refutation, or compression
measurement is claimed.
\end{document}
